\section{The Impedance of Free Space (\texorpdfstring{$Z_0$}{Z0})}
\label{sec:z0_derivation}

A foundational parameter in classical electromagnetism is the Characteristic Impedance of Free Space, $Z_0 = \sqrt{\mu_0/\epsilon_0} \approx 376.73\;\Omega$. In the VCA framework, this possesses a literal mechanical identity: it is the \textit{acoustic impedance} of the discrete lattice, encoded into an electrical constant by the Topo-Kinematic mapping.

\subsection{Derivation from the Discrete LC Ladder}

A single lattice cell of pitch $\ell_{node}$ carries distributed inductance and capacitance per unit length equal to the vacuum permeability and permittivity:
\begin{resultbox}{Per-Cell Lumped Elements}
\begin{equation}
    L_{cell} = \mu_0\, \ell_{node}, \qquad C_{cell} = \epsilon_0\, \ell_{node}
    \label{eq:cell_elements}
\end{equation}
\end{resultbox}

Evaluating numerically from the physics engine constants ($\ell_{node} = 3.862 \times 10^{-13}$ m):
\begin{align}
    L_{cell} &= 1.257 \times 10^{-6} \times 3.862 \times 10^{-13} = 4.855 \times 10^{-19} \text{ H} \\
    C_{cell} &= 8.854 \times 10^{-12} \times 3.862 \times 10^{-13} = 3.419 \times 10^{-24} \text{ F}
\end{align}
The characteristic impedance of each lattice cell is:
\begin{resultbox}{Scale-Invariant Characteristic Impedance}
\begin{equation}
    Z_{cell} = \sqrt{\frac{L_{cell}}{C_{cell}}} = \sqrt{\frac{\mu_0\, \ell_{node}}{\epsilon_0\, \ell_{node}}} = \sqrt{\frac{\mu_0}{\epsilon_0}} \equiv Z_0 \approx 376.73\;\Omega
    \label{eq:z0_cell}
\end{equation}
\end{resultbox}

The lattice pitch cancels identically. This is the fundamental reason $Z_0$ is a universal constant: it is a property of the node-to-node impedance ratio of the lattice, independent of the absolute scale $\ell_{node}$. Every cell, at every location in the universe, presents the same $376.73\;\Omega$ characteristic impedance.

\subsection{Signal Propagation Velocity}

The group velocity of a signal through a lumped LC ladder is:
\begin{resultbox}{Propagation Velocity from Discrete Components}
\begin{equation}
    v_g = \frac{\ell_{node}}{\sqrt{L_{cell}\, C_{cell}}} = \frac{\ell_{node}}{\sqrt{\mu_0\, \epsilon_0\, \ell_{node}^2}} = \frac{1}{\sqrt{\mu_0\, \epsilon_0}} \equiv c
    \label{eq:c_from_lc}
\end{equation}
\end{resultbox}

The invariant speed of light is structurally identical to the slew rate of a discrete LC transmission line. No continuous medium is assumed; $c$ emerges from lumped elements.

\subsection{Mechanical Acoustic Impedance}

Applying the Topo-Kinematic identity (Section~\ref{sec:topo_kinematic}), the mechanical impedance equivalent is:
\begin{resultbox}{Mechanical Acoustic Impedance of the Vacuum}
\begin{equation}
    Z_{mech} = \xi_{topo}^{2}\, Z_0 = (4.149 \times 10^{-7})^2 \times 376.73 \approx 6.485 \times 10^{-11} \text{ kg/s}
    \label{eq:z_mech}
\end{equation}
\end{resultbox}

This value represents the absolute force-per-velocity ratio of a single lattice node. It is the fundamental quantum of mechanical impedance in the $\mathcal{M}_A$ condensate.

\subsection{Impedance Across Physical Regimes}

Table~\ref{tab:impedance_regimes} summarises how $Z_0$ behaves across the three operating regimes of the vacuum, as established in Volume~I.

\begin{table}[H]
\centering
\caption{Characteristic impedance across the three vacuum operating regimes. Gravity preserves $Z_0$ (stealth); particle cores and event horizons destroy it ($Z \to 0$).}
\label{tab:impedance_regimes}
\small
\begin{tabular}{@{}l c c c c@{}}
\toprule
\textbf{Regime} & $\mu_{eff}$ & $\epsilon_{eff}$ & $Z = \sqrt{\mu_{eff}/\epsilon_{eff}}$ & $\Gamma$ \\
\midrule
Linear vacuum             & $\mu_0$                & $\epsilon_0$                & $Z_0 = 376.73\;\Omega$  & $0$ \\
Gravity well ($n \gg 1$)  & $n\, \mu_0$            & $n\, \epsilon_0$            & $Z_0 = 376.73\;\Omega$  & $0$ \\
Particle core ($\Delta\phi \to \alpha$) & $\to 0$ (Meissner) & $\to 0$ (dielectric collapse) & $\to 0\;\Omega$ & $-1$ \\
Event horizon             & $\to 0$ (saturated)    & $\to 0$ (saturated)         & $\to 0\;\Omega$          & $-1$ \\
\bottomrule
\end{tabular}
\end{table}

\noindent The critical distinction: gravity scales $\mu$ and $\epsilon$ \textit{symmetrically} ($n \times n$), preserving $Z_0$ and producing zero reflection. Topological saturation (particles, event horizons) drives both to zero \textit{asymmetrically} via Axiom~4, collapsing $Z$ and creating perfect mirrors ($\Gamma = -1$). This is why gravity wells are RF-transparent (stealth) while particle boundaries and event horizons are perfect reflectors.